

%=====================================================================================================================================================================================================================







\section{Derivation of the explicit form of the Uhlmann partial isometry}
\label{sec:analyze partial iso}


Following the approach in Ref.~\cite{Jozsa1994FidelityforMixedQuantumStates}, we show that the isometry $V^\sB = \sgn^{(\rm SV)}\big(\tr_{\hat{\sA}}\big[\ketbra{\sigma}{\rho}^{\hat{\sA}\sB}\big]\big)$, satisfies that $\rF(\rho^\sA, \sigma^\sA) = \rF(V^\sB\ket{\rho}^{\sA\sB}, \ket{\sigma}^{\sA\sB})$, where the sign function acts on the singular values of the input matrix, as shown in Eq.~\eqref{inteq:158}.
This implies that $V^\sB$ is the Uhlmann partial isometry.

To this end, we should recall the variational characterization of the trace norm: for any square matrix $M$, it holds that $\|M\|_1 = \max_{U}|\tr[UM]|$, where maximization is taken over all unitaries. The maximization is attained by the inverse of the unitary polar factor of $M$.
On the block specified by the left and right singular spaces corresponding to the non-zero singular values of $M$, the unitary polar factor of $M$ is uniquely determined and given by $\sgn^{(\rm SV)}(M)$. Outside the block, it can act arbitrarily, but this arbitrariness does not affect the result.


Using the above fact, we can derive the explicit form of the Uhlmann partial isometry in a straightforward manner.
By the Uhlmann's theorem, $\rF(\rho^\sA, \sigma^\sA)$ is rephrased as
\begin{align}
    \rF(\rho^\sA, \sigma^\sA) 
    &= \max_{U^\sB}\rF(U^\sB\ket{\rho}^{\sA\sB}, \ket{\sigma}^{\sA\sB}) \\
    &= \max_{U^\sB} \big|\bra{\sigma}^{\sA\sB}U^\sB\ket{\rho}^{\sA\sB}\big|^2 \\
    &= \max_{U^\sB} \big|\tr\big[U^\sB\tr_\sA[\ketbra{\rho}{\sigma}^{\sA\sB}]\big]\big|^2,
\end{align}
where maximization is taken over all unitaries $U^\sB$.
On the block specified by the left and right singular spaces corresponding to the non-zero singular values of $\tr_\sA[\ketbra{\rho}{\sigma}^{\sA\sB}]$, the unitary that achieves the maximization is uniquely determined as $\big(\sgn^{(\rm SV)}(\tr_\sA[\ketbra{\rho}{\sigma}^{\sA\sB}])\big)^\dag = \sgn^{(\rm SV)}(\tr_\sA[\ketbra{\sigma}{\rho}^{\sA\sB}])$, which is a partial isometry. We denote the partial isometry by $V^\sB$ and refer to it as the Uhlmann partial isometry.


Finally, we see that the singular values of $\tr_\sA[\ketbra{\sigma}{\rho}^{\sA\sB}]$ correspond to those of $\sqrt{\sigma^\sA}\sqrt{\rho^\sA}$.
Without loss of generality, we suppose $d_\sA \leq d_\sB$, which is justified by the fact that we can arbitrarily pad $\sB$ with the state $\ket{0}$.

Suppose that the Schmidt decomposition of the purified state $\ket{\rho}^{\sA\sB}$ and $\ket{\sigma}^{\sA\sB}$ is given by
\begin{align}
    \label{inteq:schmidt decomp rho_sigma}
      &\ket{\rho}^{\sA\sB} = \sum_{i=1}^{r_\rho}\sqrt{p_i}\ket{e_i}^\sA\ket{f_i}^\sB, \ \ \text{and} \ \ \ \ket{\sigma}^{\sA\sB} = \sum_{j=1}^{r_\sigma}\sqrt{q_j}\ket{g_j}^\sA\ket{h_j}^\sB,
\end{align}
respectively.
Then, we have that
\begin{align}
    \tr_\sA\big[\ketbra{\sigma}{\rho}^{\sA\sB}\big] 
    &= \sum_{i,j}\sqrt{p_i}\sqrt{q_j}\braket{e_i}{g_j}\ketbra{h_j}{f_i}^\sB \\
    &= \Big(\sum_{i,j}\sqrt{p_i}\sqrt{q_j}\braket{g_j}{e_i}\ketbra{h_j^*}{f_i^*}^\sB\Big)^* \\
    &= \big(V_2^{\sA\rarr\sB}\sum_{i,j}\sqrt{p_i}\sqrt{q_j}\braket{g_j}{e_i} \ketbra{g_j}{e_i}^\sA(V_1^{\sA\rarr\sB})^\dag\big)^* \\
    &= \big(V_2^{\sA\rarr\sB}\sqrt{\sigma^\sA}\sqrt{\rho^\sA}(V_1^{\sA\rarr\sB})^\dag\big)^*,
\end{align}
where $V_1^{\sA\rarr\sB}$ and $V_2^{\sA\rarr\sB}$ are isometries that map $\ket{e_i}^\sA$ to $\ket{f_i^*}^\sB$ and $\ket{g_j}^\sA$ to $\ket{h_j^*}^\sB$, respectively.


Since singular values are invariant under complex conjugation and isometries, we conclude that $\tr_\sA\big[\ketbra{\sigma}{\rho}^{\sA\sB}\big]$ and $\sqrt{\sigma^\sA}\sqrt{\rho^\sA}$ have identical singular values.
Note that all of these singular values lie within the range $[0,1]$.
The rank of $\tr_{\sA}\big[\ketbra{\sigma}{\rho}^{\sA\sB}\big]$ also coincides with the rank of $\sqrt{\sigma^\sA}\sqrt{\rho^\sA}$, as they have the same number of singular values.












%=======================================================================================================================================================






\section{Proof of \texorpdfstring{Lemma~\ref{lem:state error not accumulate}}{Lemma~\ref{lem:state error not accumulate}}}


\label{sec:appendix error not accumulate}





We now prove \Cref{lem:state error not accumulate}, which is restated below.

\lemrobstexp*


\begin{proof}[Proof of \Cref{lem:state error not accumulate}]

Observe that
\begin{align}
    \f{d}{ds}(e^{-isB}e^{isA}) 
    &= -iBe^{-isB}e^{isA} + e^{-isB}iAe^{isA} \\
    \label{inteq:91}
    &=ie^{-isB}(A-B)e^{isA}.
\end{align}
Suppose $t \geq 0$, we integrate both sides of Eq.~\eqref{inteq:91} from $0$ to $t$, obtaining that
\begin{align}
\label{inteq:56}
    e^{-itB}e^{itA} - \bI = \int_0^t ie^{-isB}(A-B)e^{isA} ds.
\end{align}
Multiplying both sides of Eq.~\eqref{inteq:56} from the left by $e^{itB}$ and taking the operator norm gives
\begin{align}
    \big\|e^{itA} - e^{itB}\big\|_\infty 
    &= \Big\|\int_0^t ie^{i(t-s)B}(A-B)e^{isA} ds\Big\|_\infty \\
    &\leq \int_0^t |i|\big\|e^{i(t-s)B}\big\|_\infty \|A-B\|_\infty\big\|e^{isA}\big\|_\infty ds \\
    &= \int_0^t ds \|A-B\|_\infty \\
    &= t \|A-B\|_\infty.
\end{align}
In the second equation, we used that $e^{i(t-s)B}$ and $e^{isA}$ are unitaries, since $A$ and $B$ are Hermitian.

When $t < 0$, a similar approach shows that $\big\|e^{itA} - e^{itB}\big\|_\infty \leq -t\|A-B\|_\infty$.
Hence, $\big\|e^{itA} - e^{itB}\big\|_\infty \leq |t|\|A-B\|_\infty$ holds for any $t \in \bR$, and the proof is completed.



\end{proof}





%================================================================================================================================================


